Aquí desarrollaremos nuestro trabajo. Contrastaremos la ideas entre varios autores y, si es posible, con las nuestras. Podemos incluir las secciones y subsecciones que necesitemos.
\newline

Tanto en esta Sección, como en el resto del Trabajo, se debe cuidar la redacción, el uso de las comas, puntos, puntos y coma, puntos y aparte y demás recursos. Fijaos en esta plantilla como estan escritas ciertas palabras que usamos mucho: \emph{Sección xx}, pero \emph{secciones}; \emph{Figura xx}, pero \emph{figuras}; \emph{Ecuación xx}, pero \emph{ecuaciones}; etc... Podéis echar un vistazo a trabajos de fin de estudios para comprobar la organización, ortografía y redacción de un trabajo a este nivel. En el repositorio de UNIR tenéis la posibilidad de acceder a trabajos presentados como TFMs en esta Titulación que os servirán de modelo: \url{https://reunir.unir.net/handle/123456789/6218}
\newline

Las figuras y tablas tienen que tener su numeración y su título. Además, se debe indicar la fuente, si es de elaboración propia o se ha tomado de alguna referencia.
\newline

Por ejemplo: